\subsection{Homotopy-coherent inputs}
\label{subsec:2dloc-inputs}

We list the results that produce or compare homotopy-coherent data. Some of them are proved in other models of $(\infty,2)$-categories.

\subsubsection*{(I0) Comparison of models}
The $\infty$-category $\twoCat$ is equivalent to the $\infty$-categories underlying the other standard models, e.g.\ complete $2$-fold Segal spaces \cite{gepnerhaugseng2015enriched} and categories strictly enriched in quasi-categories \cite{haugseng2015rectification}. By the unicity theorem of Barwick--Schommer-Pries \cite{barwickschommerpries2021unicity}, any two equivalences between such $\infty$-categories differ by an automorphism generated by $(-)\opcat$ and $(-)^{\mathrm{co}}$.
The statements (I1)--(I2) below, taken together with their duals, are invariant under these automorphisms, so it does not matter in which model they are proved.
\todo{Check the precise statements of the comparison theorems in \cite{gepnerhaugseng2015enriched,haugseng2015rectification}.}

\subsubsection*{(I1) The walking adjunction}
Let $\Adj$ be the $2$-category freely generated by an adjunction $l \dashv r$ between two objects $-$ and $+$, viewed as an $(\infty,2)$-category, and let $\ell,\rho\colon \Delta^{1} \to \Adj$ select $l\colon - \to +$ and $r\colon + \to -$.

\begin{thm}[Riehl--Verity; Abell\'an]
\label{thm:walking-adjunction}
For every $(\infty,2)$-category $\A$, restriction along $\ell$
\[
\ell^{*}\colon \mapp_{\twoCat}(\Adj,\A) \longrightarrow \mapp_{\twoCat}(\Delta^{1},\A)
\]
is a monomorphism of spaces whose image consists of the components of the left adjoints (in the sense of \cref{def:adjunction}). Dually, $\rho^{*}$ is a monomorphism onto the right adjoints.
\end{thm}

This is the only place where an adjunction in $\htwo\A$ is lifted to homotopy-coherent data. Where it is proved:
\begin{itemize}
\item Riehl--Verity \cite{riehlverity2016adjunctions} work with a category $\mathcal{K}$ strictly enriched in quasi-categories. They show that $\Adj$ is a cofibrant simplicial category (a simplicial computad), that every adjunction in the homotopy $2$-category of $\mathcal{K}$ extends to a simplicial functor $\Adj \to \mathcal{K}$, and that the relevant spaces of extensions are contractible Kan complexes \cite[\S 4]{riehlverity2016adjunctions}.
To deduce \cref{thm:walking-adjunction} from this one needs, in addition, that for cofibrant $\Adj$ and fibrant $\mathcal{K}$ these spaces of simplicial functors compute mapping spaces in $\twoCat$; this is where (I0), in the form of \cite{haugseng2015rectification}, enters.
\todo{Check which of the extension statements of \cite[\S 4]{riehlverity2016adjunctions} starts from the left adjoint alone, and whether the passage to $\twoCat$ is written out in \cite{haugseng2021laxmonads}.}
\item Abell\'an \cite{abellan2026bifibrations} gives a model-independent proof of the universal property of $\Adj$, as an application of free bifibrations of $(\infty,2)$-categories. \todo{Theorem number.}
\end{itemize}

\subsubsection*{(I2) Adjoints in functor $(\infty,2)$-categories}
\begin{thm}[Abell\'an--Gagna--Haugseng]
\label{thm:adjoints-in-functor-cats}
Let $\A$, $\B$ be $(\infty,2)$-categories and $\alpha\colon F \Rightarrow G$ a $1$-morphism of $\func(\A,\B)$. Then $\alpha$ is a left adjoint in $\func(\A,\B)$ if and only if $\alpha_{a}$ is a left adjoint in $\B$ for every object $a$ and, for every $1$-morphism $t\colon a \to a'$ of $\A$, the lax square
\[
\begin{tikzcd}
F(a) \ar[r, "F(t)"] \ar[d, "\alpha_{a}"'] & F(a') \ar[d, "\alpha_{a'}"] \ar[dl, Rightarrow, shorten=3mm, "\alpha_{t}"'] \\
G(a) \ar[r, "G(t)"'] & G(a')
\end{tikzcd}
\]
(the naturality square of $\alpha$ at $t$, with $\alpha_{t}\colon \alpha_{a'}F(t) \simeq G(t)\alpha_{a}$) is right Beck--Chevalley.
Dually, $\alpha$ is a right adjoint if and only if every $\alpha_{a}$ is a right adjoint and every transposed naturality square, with $1$-morphisms $\alpha_{a}$, $\alpha_{a'}$ horizontal and $2$-morphism $\alpha_{t}^{-1}\colon G(t)\alpha_{a} \Rightarrow \alpha_{a'}F(t)$, is left Beck--Chevalley.
\end{thm}

This is proved in \cite{abellangagnahaugseng2024straightening}, as a special case of their characterization of adjoints in $(\infty,2)$-categories of functors and (op)lax transformations; the dual statement follows by applying $(-)^{\mathrm{co}}$.
\todo{Theorem number, and check that the strong case is stated in exactly this form.}
Note that \cref{thm:adjoints-in-functor-cats} is \emph{not} a formal consequence of the corresponding statement for $2$-categories: $\htwo\func(\A,\B)$ is in general not the $2$-category of pseudofunctors $\htwo\A \to \htwo\B$, because homotopy categories do not commute with the limits computing the hom-categories of $\func(\A,\B)$.

\subsubsection*{(I3) Enriched Yoneda lemma and free cocompletion}
We use the enriched Yoneda lemma and the universal property of enriched presheaf categories of Hinich \cite{hinich2020yoneda} and Heine \cite{heine2023enriched}, which are formulated in the enriched model.

\subsubsection*{(I4) Enriched adjunctions and monads}
Let $L\colon \D \to \E$ be a $\Cat$-linear left adjoint between presentable $\infty$-categories with $\Cat$-action (see below), with right adjoint $\iota$. Then $\iota$ is canonically lax $\Cat$-linear \cite[7.3.2.7]{ha}, hence an enriched functor, and $L \dashv \iota$ is an adjunction in the $(\infty,2)$-category of large $(\infty,2)$-categories, with unit and counit enriched natural transformations \cite{heine2023enriched}.
Consequently $T := \iota L$ is an enriched functor with enriched natural transformations $\eta\colon \mathrm{id} \Rightarrow T$ and $\mu := \iota\varepsilon L\colon TT \Rightarrow T$, satisfying the monad identities in $\htwo$. On underlying $\infty$-categories $T$ is an associative algebra in $\func(\iota_{1}\D,\iota_{1}\D)$ \cite[\S 4.7.3]{ha}; monads in an $(\infty,2)$-category are compared with such algebras in \cite{haugseng2021laxmonads}.
\todo{Check that the passage from lax-linear functors and transformations to enriched ones in \cite{heine2023enriched} is compatible with natural transformations.}

\subsubsection*{Presentable $(\infty,2)$-categories}
We take a \emph{presentable $(\infty,2)$-category} to be a $\Cat$-module in $\prl$, and set $\tprl := \modcat{\Cat}{\prl}$, the $\infty$-category of presentable $\infty$-categories with a colimit-preserving action of $\Cat$ (for the cartesian product) and $\Cat$-linear left adjoints.
By \cite{gepnerhaugseng2015enriched,heine2023enriched} every $\D \in \tprl$ is canonically $\Cat$-enriched, with $\func(K,\Hom_{\D}(d,d')) \simeq \mapp_{\D}(K \otimes d,d')$, tensored and cotensored over $\Cat$, and every $\Cat$-linear functor is canonically enriched.
\todo{Precise references.}
For a small $(\infty,2)$-category $\A$ the enriched presheaf category $\Free(\A) := \func(\A\opcat,\Cat)$ is in $\tprl$, and by (I3) restriction along the Yoneda embedding $y\colon \A \to \Free(\A)$ is an equivalence
\begin{equation}
\label{eq:free-univ}
\func^{L}_{\Cat}\bigl(\Free(\A),\D\bigr) \xrightarrow{\ \simeq\ } \func(\A,\D)
\end{equation}
for every $\D \in \tprl$, where the left-hand side is the $(\infty,2)$-category of $\Cat$-linear left adjoints.
A functor $f\colon \A \to \B$ induces $\Free(f) = f_{!}\colon \Free(\A) \to \Free(\B)$ (enriched left Kan extension) with right adjoint the restriction $f^{*}$, and $\Free\colon \twoCat \to \tprl$ preserves colimits.
We write $\Edge := \Free(\Delta^{1})$. By \eqref{eq:free-univ}, a $1$-morphism $s\colon x \to y$ of $\D \in \tprl$ is the same as a map $s\colon \Edge \to \D$ in $\tprl$, and by \cref{thm:walking-adjunction} a left adjoint $1$-morphism of $\D$ extends essentially uniquely to a map $\Free(\Adj) \to \D$ along $\Free(\ell)$.
